\chapter{Model Predictive Control Framework}\label{ch:mpc_framework}

The models developed in Chapter~\ref{ch:data_driven_predictor_design} predict the response of measured neural activity to a given stimulation sequence.
To use these models for control, the stimulation sequence must be chosen to reduce seizure-related activity while limiting stimulation effort and satisfying current constraints.
This chapter formulates that choice as a finite-horizon optimal control problem.
Repeatedly solving this problem from updated measurements and applying only the first input constitutes \ac{mpc}.

The chapter first establishes the general \ac{mpc} formulation and compares numerical solution methods.
It then specifies the predictor dynamics, current constraints and control costs for the neurostimulation \ac{ocp}.
Cost and constraint formulations that promote sparse stimulation are developed next.
The chapter ends with a discussion of limitations of control with learned predictors.

\section{General \ac{mpc} Formulation}\label{sec:mpc_formulation}

We consider the discrete-time prediction model
%
\begin{equation}
    \mathbf{x}_{k+1} = \mathbf{f}(\mathbf{x}_k,\mathbf{u}_k)\,,
    \qquad \mathbf{f}:\mathbb{R}^{\ndim{x}}\times\mathbb{R}^{\ndim{u}}\to\mathbb{R}^{\ndim{x}}\,,
    \label{eq:mpc_prediction_dynamics}
\end{equation}
%
where $\mathbf{x}_k\in\mathbb{R}^{\ndim{x}}$ is the model state and $\mathbf{u}_k\in\mathbb{R}^{\ndim{u}}$ is the input at time step $k\in\mathbb{N}_0$.
At each update, measurements and past inputs provide the initial state $\hat{\mathbf{x}}_k$ from which future trajectories are predicted.
For the predictors of Chapter~\ref{ch:data_driven_predictor_design}, this state contains the observation and input histories defined there.

Starting from $\mathbf{x}_{0|k}=\hat{\mathbf{x}}_k$, a candidate input sequence $\mathbf{u}_{0|k},\ldots,\mathbf{u}_{H-1|k}$ determines the predicted states through~\eqref{eq:mpc_prediction_dynamics}.
The index $j|k$ denotes a prediction $j$ steps ahead of the current update $k$.
Each input is held constant for one sampling interval $\Delta t>0$, so the horizon of $H\in\mathbb{N}_{>0}$ steps covers the duration $T_{\mathrm{pred}}=H\Delta t$.

\begin{samepage}
The input sequence is chosen to minimize a cost over this horizon while satisfying the state and input constraints.
With real-valued stage cost $\ell:\mathbb{R}^{\ndim{x}}\times\mathbb{R}^{\ndim{u}}\to\mathbb{R}$ and terminal cost $V_\mathrm{f}:\mathbb{R}^{\ndim{x}}\to\mathbb{R}$, the resulting \ac{ocp} takes the standard form~\cite{Rawlings2017}
%
\begin{equation}
    \begin{aligned}
        \min_{\mathbf{u}_{0|k},\ldots,\mathbf{u}_{H-1|k}}
        \quad & \mathcal{J}
        = \sum_{j=0}^{H-1}
          \ell(\mathbf{x}_{j|k},\mathbf{u}_{j|k})
          + V_\mathrm{f}(\mathbf{x}_{H|k}) \\
        \text{s.t.}\quad
        & \mathbf{x}_{0|k}=\hat{\mathbf{x}}_k, \\
        & \mathbf{x}_{j+1|k}=\mathbf{f}(\mathbf{x}_{j|k},\mathbf{u}_{j|k}),
          \qquad j=0,\ldots,H-1, \\
        & \mathbf{x}_{j|k}\in\mathbb{X},
          \qquad j=0,\ldots,H, \\
        & \mathbf{u}_{j|k}\in\mathbb{U},
          \qquad j=0,\ldots,H-1, \\
        & \mathbf{x}_{H|k}\in\mathbb{X}_\mathrm{f}\,.
    \end{aligned}
    \label{eq:mpc_ocp}
\end{equation}
\end{samepage}
%
Here, the predicted states are determined by the initial condition and dynamics for each candidate input sequence.
The stage cost $\ell$ can penalize deviations from a desired operating condition together with input effort.
Its weights set the relative penalties on state deviation and input effort.
Unlike the training loss $\mathcal L$, which is used to fit the predictor, this cost is minimized over future inputs with the model held fixed.

The sets $\mathbb{X}$ and $\mathbb{U}$ contain the admissible states and inputs.
Membership in these sets is required for a sequence to be feasible, regardless of its cost.
If no state restrictions are imposed, $\mathbb{X}=\mathbb{R}^{\ndim{x}}$.
In particular, penalizing seizure-related activity in $\ell$ allows its reduction to be traded against stimulation effort without requiring a hard bound on the predicted activity.  % this is a general formulation. Any seizure related talk is displace. This whole sentence/paragraph doesn't make sense: "In particular" what?

Solving~\eqref{eq:mpc_ocp} yields an input sequence $\mathbf{u}^{\star}_{0|k},\ldots,\mathbf{u}^{\star}_{H-1|k}$.
The controller applies its first element,
%
\begin{equation*}
    \mathbf{u}_k=\mathbf{u}^{\star}_{0|k}\,.
\end{equation*}
%
At time step $k+1$, new measurements provide the next initial state, and the problem is solved again with the horizon shifted forward by one step.
The remaining inputs of the previous solution are therefore replaced by a new plan based on the measured response.
Although each optimization computes an open-loop input sequence, replanning from updated measurements introduces feedback.
Figure~\ref{fig:receding_horizon} shows how the plans change as new measurements become available.

\begin{figure}
    \centering
    \includegraphics{figures/05_mpc_framework/receding_horizon}
    \caption{Illustration of receding-horizon control for a scalar system.
    Dashed curves show selected future state and input plans; solid curves show the realized trajectory and applied inputs.
    Replanning occurs at every sampling instant, with plans at $k=0$, $3$ and $6$ shown for clarity.
    The realized state departs from the nominal forecasts because of disturbances.
    The dotted horizontal line marks the upper input bound.}\label{fig:receding_horizon}
\end{figure}

Finite-horizon optimization alone does not guarantee closed-loop stability or recursive feasibility. % this kind of falls out of the sky. What even is recursive feasibility? 
The terminal cost $V_\mathrm{f}$ can approximate the remaining cost of continuing from $\mathbf{x}_{H|k}$. % what is remaining cost? we need to explain somewhere that the finite horizon form is an approximation of the infinite horizon problem
The terminal set $\mathbb{X}_\mathrm{f}\subseteq\mathbb{X}$ can restrict this endpoint to states from which admissible operation can continue.
Setting $\mathbb{X}_\mathrm{f}=\mathbb{X}$ removes the additional endpoint restriction.
When a local controller preserves the terminal set and the terminal cost satisfies suitable decrease conditions, these ingredients can establish stability and recursive feasibility~\cite{Mayne2000}.
Recursive feasibility means that initial feasibility is preserved at subsequent updates.
Merely assigning a cost to the final prediction does not establish either property.

The \ac{ocp} defines the control objective and admissible predictions.
The next section compares direct multiple and single shooting as two numerical representations of this problem.

\section{Numerical Transcription and Solution}\label{sec:mpc_numerical_solution}

The numerical solution can treat predicted states as additional optimization variables (direct multiple shooting) or compute them from the inputs (direct single shooting).

Direct multiple shooting introduces the predicted states of~\eqref{eq:mpc_ocp} as independent optimization variables alongside the inputs.
For one-step shooting intervals, it optimizes $(H+1)\ndim{x}$ state components together with $H\ndim{u}$ input components and enforces the dynamics through the continuity constraints~\cite{Bock1984}
%
\begin{equation}
    \mathbf{x}_{0|k} = \hat{\mathbf{x}}_k, \qquad
    \mathbf{x}_{j+1|k} - \mathbf{f}\!\left(\mathbf{x}_{j|k}, \mathbf{u}_{j|k}\right)
    = \mathbf{0},
    \quad j = 0, \dots, H-1 \,. \label{eq:multiple_shooting}
\end{equation}
%
At an intermediate iterate, the equalities in~\eqref{eq:multiple_shooting} may not hold.
The iterate then consists of one-step trajectory pieces with gaps between them, as shown in Figure~\ref{fig:shooting}.

Direct single shooting implements the input-only formulation of~\eqref{eq:mpc_ocp} by computing the predicted states through model rollout.
Only the $H\ndim{u}$ input components are optimized.
For a state $\mathbf{x}\in\mathbb{R}^{\ndim{x}}$ and inputs $\mathbf{u}_i\in\mathbb{R}^{\ndim{u}}$, $i=0,\ldots,H-1$, define the rollout maps $\boldsymbol{\varphi}_j$ recursively by
%
\begin{align}
    \boldsymbol{\varphi}_0(\mathbf{x})
    &= \mathbf{x}, \nonumber \\
    \boldsymbol{\varphi}_1(\mathbf{x}, \mathbf{u}_0)
    &= \mathbf{f}(\mathbf{x}, \mathbf{u}_0), \nonumber \\
    \boldsymbol{\varphi}_{j+1}(\mathbf{x}, \mathbf{u}_0, \dots, \mathbf{u}_j)
    &= \mathbf{f}\!\left(
       \boldsymbol{\varphi}_j(\mathbf{x}, \mathbf{u}_0, \dots, \mathbf{u}_{j-1}),
       \mathbf{u}_j\right),
       \qquad j = 1, \dots, H-1 \,. \label{eq:single_shooting}
\end{align}
%
By~\eqref{eq:single_shooting}, the predicted states for $j=1,\ldots,H$ satisfy $\mathbf{x}_{j|k} = \boldsymbol{\varphi}_j(\hat{\mathbf{x}}_k,
\mathbf{u}_{0|k}, \dots, \mathbf{u}_{j-1|k})$.
Substitution into the cost and state constraints leaves a numerical problem in the inputs alone, with the dynamics satisfied by construction.
Every iterate defines a dynamically consistent rollout of the dynamics model, although it does not necessarily satisfy the remaining constraints.

State constraints act directly on multiple-shooting variables.
Under single shooting, they act on the rollout maps $\boldsymbol{\varphi}_j$.
Input constraints retain their original form in either transcription.

The choice of optimization variables also affects how model sensitivities enter the numerical problem.
At predicted step $q=0,\ldots,H-1$, define the Jacobians $\mathbf{A}_q = \partial\mathbf{f}/\partial\mathbf{x}
(\mathbf{x}_{q|k},\mathbf{u}_{q|k})$ and $\mathbf{B}_q = \partial\mathbf{f}/\partial\mathbf{u}
(\mathbf{x}_{q|k},\mathbf{u}_{q|k})$,
with $\mathbf{A}_q\in\mathbb{R}^{\ndim{x}\times\ndim{x}}$ and $\mathbf{B}_q\in\mathbb{R}^{\ndim{x}\times\ndim{u}}$.
For $0\le r<j\le H$, repeated application of the chain rule gives
%
\begin{equation*}
    \frac{\partial \mathbf{x}_{j|k}}{\partial \mathbf{u}_{r|k}}
    = \mathbf{A}_{j-1}\mathbf{A}_{j-2}\cdots\mathbf{A}_{r+1}\mathbf{B}_r \,,
\end{equation*}
%
where the matrix product is the identity when $r=j-1$.
The same Jacobian product appears when an iterated prediction model is trained through a multiple-step rollout.
The product may amplify or attenuate early-input sensitivities, impairing numerical conditioning.
Single shooting therefore gives a smaller optimization problem, but its derivatives involve products across the prediction horizon.
Multiple shooting introduces additional variables and equality constraints, whose Jacobians involve local model derivatives and have a sparse structure~\cite{Bock1984,Rawlings2017}.
The computational tradeoff depends on the dynamics, horizon length and numerical solver.

\begin{figure}
    \centering
    \includegraphics{figures/05_mpc_framework/shooting}
    \caption{Direct multiple and single shooting over one prediction horizon.
    Multiple shooting treats the state at each sampling instant as an optimization variable; gaps between the one-step predictions and the next state variables indicate violated dynamics constraints at an intermediate iterate.
    Single shooting computes one continuous model rollout from the initial state, with only the inputs as optimization variables.}\label{fig:shooting}
\end{figure}

For a nonlinear dynamics model, either transcription gives a finite-dimensional \ac{nlp}.
Both can be solved using \ac{sqp}, which forms quadratic subproblems with linearized constraints~\cite{Nocedal2006}.
Alternatively, primal-dual interior-point methods, such as IPOPT, solve barrier subproblems while reducing the barrier parameter~\cite{Waechter2006,Nocedal2006}.
The following section specializes the general \ac{ocp} to the learned neurostimulation predictors, current constraints and control costs used in this thesis.

\section{\ac{mpc} with Learned Input-Output Predictors}\label{sec:mpc_input_output_formulation}

\subsection{History-State Realization}
% Refer to Chapter 4 for the observable and current histories rather than repeating their derivation.
% Distinguish the learned output map o_{j+1|k} = f_theta(h_{j|k}, u_{j|k}) from the full history-state transition.
% Construct the latter by shifting the histories and appending the predicted observable and candidate current.
% Identify x = h in the general OCP only after defining this complete transition map.
% Explain that most history-state coordinates are copied entries, not newly predicted quantities.
% Initialize from measured observations and past currents, without an observer of the biological state.
% Specify waveform-sample and spectral-frame time grids, controller updates and zero-order hold actuation.

\subsection{\ac{narx} Transcription}
% Motivate an output-lifted alternative to lifting the complete history state in multiple shooting.
% Introduce future outputs and currents as decision variables and enforce the learned output-map equalities.
% Reconstruct histories from fixed measured pre-horizon entries and the future output/current variables.
% Explain how this eliminates redundant shift-register variables and copy constraints.
% Develop the partially condensed variant: lift observation-history windows only at segment boundaries,
% roll out between boundaries and reconstruct input histories directly from the current decision sequence.
% Distinguish fully output-lifted one-step defects from boundary-window defects in the partially condensed variant.
% State decision-variable and defect counts relative to full-history multiple shooting and single shooting.
% Present algebraic structure here; runtime, convergence and the empirical choice of SS belong in Chapter 6.
% Source: feat/narx_mpc, src/neuro/control/narx_mpc.py.
% Benchmark context: knowledge-base/Notes/narx_mpc_benchmark.md.

% State the specific input constraints in a short paragraph, without a separate subsection:
% hard per-electrode bounds |u_p| <= 2 mA and Kirchhoff equality sum_p u_p = 0.
% Retain these constraints in each transcription and distinguish them from soft suppression penalties.

\subsection{Control Cost Functions}\label{sec:mpc_cost_functions}

% Motivate the use of measurable proxies: recruitment is the ultimate control objective, but its
% regional plant activity is unavailable to the controller, so the recruitment measure cannot be used directly.
% Instead, formulate suppression costs from predicted scalp measurements.

\subsubsection*{Healthy References}
% Explain the construction of empirical healthy channel means and channel-frequency log-power envelopes.
% Use the same spectral transformation for healthy references and predicted quantities.

\subsubsection*{Waveform Tracking Cost}
% Define squared waveform deviation from empirical healthy channel means.
% This tracks a healthy operating point, not the phase of a prescribed healthy waveform.
% Retain the implemented quadratic tracking cost.

\subsubsection*{Spectral Suppression Cost}
% Define the squared one-sided hinge [predicted log power - healthy envelope]_+^2.
% Explain why power below the healthy envelope is unpenalized.
% Define weighting and averaging over scored times, channels and frequency bins.
% Introduce the two evaluations of this objective:
% (1) waveform predictor with spectral cost computed from its predicted waveform continuation;
% (2) spectral predictor with the cost applied directly to predicted spectral features.
% For waveform-derived spectra, explain the measured boundary history and overlapping spectral windows.
% Overlapping STFT windows make the waveform-derived spectral cost depend jointly on multiple prediction steps, extending the stage-additive form in \eqref{eq:mpc_ocp}.
% Together with waveform tracking, these establish the three predictor-cost combinations.

\subsubsection*{Quadratic Current Effort and Combined Objective}
% Define the squared L2 current-effort term.
% Interpret it as a current-effort measure; electrical energy requires an additional voltage/impedance model.
% Assemble the suppression and stimulation terms into a modular weighted objective.
% Components may be combined as required; the experiments specify the combinations investigated.
% Explain horizon normalization and scoring of the final predicted sample or Frame.

% Collect the specific predictor, constraints and modular cost into our OCP.
% Explain final-prediction scoring and distinguish it from a stabilizing terminal construction.

\section{Sparsity-Promoting Control Formulations}\label{sec:mpc_sparse_formulations}

% Introduce L1 current penalties and Active Interval optimization as distinct sparsity mechanisms.
% These options can replace or supplement the quadratic current-effort term.
% Give the general formulation and specialize it to the spectral predictor and spectral suppression cost.
% The experimental combinations and their comparative performance belong in Chapter 6.

\subsection{L1 Current Cost and Exact Epigraph Formulation}
% Define the L1 penalty on physical electrode currents and explain its sparsity-promoting role.
% Introduce auxiliary variables with -z <= u <= z and minimize their weighted sum.
% Develop the exact epigraph formulation; do not develop the smooth approximation as a study method.
% Distinguish summed absolute current from signed net charge and from the number of enabled intervals.

\subsection{Active Interval Optimization}
% Introduce one shared binary variable delta_{j|k} per horizon interval.
% Gate all physical electrode currents through -u_max,p delta <= u_p <= u_max,p delta.
% Retain Kirchhoff's equality and the original predictor dynamics driven by physical currents.
% Define both alternatives: a normalized activation penalty and a hard horizon-local activation cap.
% Explain that the implementation treats penalty and cap as alternative options.
% An enabled interval permits stimulation but need not contain nonzero current.
% Count enabled intervals, not pulse onsets; adjacent enabled intervals may have different currents.

\subsection{Mixed-Integer Nonlinear Programming with Bonmin}
% Formulate the binary-activation problem as a MINLP and explain the relevant solution principles.
% Explain the implemented Bonmin solution route without claiming global optimality for nonconvex problems.
% Report observed computational limitations only in Chapter 6.

\subsection{Continuous Relaxation and Homotopy}
% Relax activation variables to [0, 1] and add beta sum_j delta_j(1-delta_j).
% Explain continuation from the unpenalized relaxation through increasing integrality-penalty weights.
% Describe warm starts between continuation passes, integrality tolerance and termination.
% The implementation retains sufficiently integral solutions without a separate final rounding/recovery solve.
% Explain that finite local continuation does not guarantee a globally optimal binary solution.

\subsection{Candidate \ac{mpc} with Grouped Activation Decisions}
% Group neighboring horizon intervals into blocks sharing an activation decision.
% Enumerate the 2^B block schedules and solve the continuous-current NLP for each schedule.
% Fix currents to zero in disabled intervals; keep currents independently optimized within enabled blocks.
% Add the enabled-interval penalty and select the lowest-cost successful candidate.
% Unequal block lengths contribute according to enabled interval count, not enabled block count.
% Explain that the implemented Candidate MPC uses the penalty formulation, without an explicit activation cap.
% Exhaustive enumeration covers the restricted schedule set; local NLP solves do not establish global optimality.

\section{Limitations of \ac{mpc} with Learned Predictors}\label{sec:mpc_learned_predictor_limitations}

% Explain how optimization can exploit forecast errors and inaccurate learned current sensitivities.
% Discuss distribution shift under optimized stimulation relative to identification data.
% Distinguish enforcing current constraints from guaranteeing seizure suppression or biological safety.
% State the limits of the adopted nominal formulation and numerical solution guarantees.
% Connect these methodological limitations to the evaluation questions in Chapter 6 without anticipating results.
